%% =====================================================================
%%  Corrigés — Unit 10 « How much is it? »
%% =====================================================================
\corrigesUnite{10}{How much is it?}

\begin{corrige}{1}
\begin{reponses}[4]
\item \rep{a sweater}
\item \rep{a dress}
\item \rep{a coat}
\item \rep{shoes}
\item \rep{a scarf}
\item \rep{pants}
\item \rep{a skirt}
\item \rep{gloves}
\end{reponses}
\retiens{\en{pants}, \en{jeans}, \en{shoes}, \en{gloves}, \en{glasses} : toujours au pluriel.}
\end{corrige}

\begin{corrige}{2}
\begin{reponses}[4]
\item \rep{an}
\item \rep{some}
\item \rep{some}
\item \rep{an}
\item \rep{some}
\item \rep{some}
\item \rep{a}
\item \rep{some}
\end{reponses}
\pourquoi{\en{milk}, \en{bread}, \en{coffee} et \en{money} sont indénombrables : jamais de \en{a}.
\en{eggs} est un pluriel : \en{some}.}
\end{corrige}

\begin{corrige}{3}
\begin{reponses}[3]
\item \rep{How much}
\item \rep{How many}
\item \rep{How much}
\item \rep{How much}
\item \rep{How many}
\item \rep{How much}
\end{reponses}
\pourquoi{\en{how much} pour un prix (1 et 4) ou un indénombrable (\en{milk}, \en{money}) ; \en{how
many} devant un pluriel (\en{eggs}, \en{apples}).}
\end{corrige}

\begin{corrige}{4}
\begin{reponses}[3]
\item \rep{some}
\item \rep{any}
\item \rep{any}
\item \rep{some}
\item \rep{any}
\item \rep{some}
\end{reponses}
\pourquoi{phrase 4 : on demande quelque chose poliment, \en{some} reste ; phrases 2 et 5 :
négations, donc \en{any}.}
\end{corrige}

\begin{corrige}{5}
\begin{reponses}[2]
\item \rep{twelve dollars}
\item \rep{thirty dollars}
\item \rep{ninety-nine dollars}
\item \rep{seventy-five cents}
\item \rep{fifteen fifty}
\item \rep{two ninety-nine}
\end{reponses}
\pourquoi{pour un prix avec des cents, on dit le plus souvent les deux nombres à la suite :
\en{fifteen fifty} ; la version complète est \en{fifteen dollars and fifty cents}.}
\end{corrige}

\begin{corrige}{6}
\begin{reponses}[5][french]
\item \rep{b}
\item \rep{c}
\item \rep{a}
\item \rep{e}
\item \rep{d}
\end{reponses}
\end{corrige}

\begin{corrige}{7}
\begin{reponses}[4]
\item \rep{30}
\item \rep{14}
\item \rep{50}
\item \rep{16}
\item \rep{70}
\item \rep{18}
\item \rep{90}
\item \rep{\$15}
\end{reponses}
\pourquoi{écoutez où tombe l'accent et si le \phon{n} final est présent. Texte de la piste
page~\pageref{trans-10.10}.}
\end{corrige}

\begin{corrige}{8}
\begin{reponses}[4]
\item \rep{\$45}
\item \rep{\$80}
\item \rep{\$13}
\item \rep{\$99}
\end{reponses}
\pourquoi{le vendeur précise \en{Not thirty: thirteen!} pour éviter toute confusion avec \en{thirty}.
Le manteau est en promotion (\en{on sale}).}
\end{corrige}

\begin{corrige}{9}
\begin{reponses}[4]
\item \rep{any}
\item \rep{isn't}
\item \rep{some}
\item \rep{many}
\item \rep{some}
\item \rep{much}
\item \rep{dollars}
\end{reponses}
\end{corrige}

\begin{corrige}{10}
\begin{reponses}[1]
\item \rep{How much is this coat? — It's a hundred dollars.}
\item \rep{Can I try it on?}
\item \rep{There isn't any milk.} \textit{ou} \rep{There's no milk.}
\item \rep{How many eggs are there?}
\item \rep{We need some bread and some cheese.}
\item \rep{These shoes are expensive.}
\item \rep{I'll take it. Can I pay cash?}
\item \rep{I drink a lot of coffee.}
\end{reponses}
\pourquoi{phrase 6 : \en{shoes} est pluriel, donc \en{These… are}. Phrase 7 : \en{pay cash} = payer
en espèces (\en{pay by card} = par carte).}
\end{corrige}

\begin{corrige}{11}
Réponse modèle :
\begin{anglais}\otTableau
\begin{tabular}{@{}l@{\hspace{6mm}}l@{}}
\otNomDit{vendeur} & Hi! Can I help you?\\
\otNomDit{vous}    & Yes, please. How much are these jeans?\\
\otNomDit{vendeur} & They're forty-nine dollars.\\
\otNomDit{vous}    & Can I try them on?\\
\otNomDit{vendeur} & Sure. The fitting room is over there.\\
\otNomDit{vous}    & Great, I'll take them. Can I pay by card?\\
\end{tabular}
\end{anglais}
\retiens{pour un objet au pluriel (\en{jeans}, \en{shoes}), on dit \en{them} : \en{Can I try them on?},
\en{I'll take them.}}
\end{corrige}

\begin{corrige}{12}
Réponses modèles :
\begin{reponses}[1]
\item \rep{Yes, please. I'm looking for a sweater. How much is this one?}
\item \rep{Can I try it on?}
\item \rep{Medium, please.}
\item \rep{Yes, I love it! I'll take it.}
\item \rep{By card, please.} \textit{ou} \rep{Cash.}
\end{reponses}
\retiens{\en{I'm looking for…} = « je cherche… » ; \en{this one} = « celui-ci ».}
\end{corrige}

\begin{corrigetest}
\begin{reponses}[5][french]
\item \rep{b}
\item \rep{c}
\item \rep{a}
\item \rep{b}
\item \rep{a}
\item \rep{c}
\item \rep{b}
\item \rep{a}
\item \rep{c}
\item \rep{b}
\end{reponses}
\pourquoi{question 5 : \en{bread} est indénombrable : \en{some bread}, jamais \en{a bread}.}
\end{corrigetest}
